\section{The charge algebra, domains, and rotational invariance}
\label{sec:algebra}

We establish the operator framework needed for the index and limiting
arguments: closed charge domains, a charge coercive on a chosen space
of torus invariants, and even rotationally invariant undeformed kernel
vectors.

We use the real Euclidean space \(\mathfrak k\) of traceless Hermitian
matrices, with scalar product \(\Tr(XY)\) and real Lie bracket
\([X,Y]=(XY-YX)/\ii\). Thus \(\ad(X)\) is a real skew matrix on
\(\mathfrak k\). Uppercase Roman indices denote an orthonormal color
basis and Greek indices denote spinor components;
\(i,j=1,\ldots,9\), \(a,b,c=1,2,3\), and \(d=4,\ldots,9\) are spatial
indices. For \(w\in\R^{16}\) and \(z\in\mathfrak k\), write
\(\theta(w)\cdot z=\sum_{\alpha,A}\theta_\alpha^A w_\alpha z^A\).
All Clifford tensor products are graded tensor products.

Set
\[
 T=\gamma^1\gamma^2\gamma^3,\qquad
 s_a=2,\quad s_d=-1,
 \qquad A_i(v)=\gamma^i v,\quad
 B_{ij}(v)=-\gamma^{ij}v,\quad L_i(v)=s_i\gamma^iTv.
\]
Here \(T^t=-T\), \(T^2=-\Id\), and \(T\) commutes with the first
three gamma matrices and anticommutes with the other six. For a unit
real spinor \(v\), coupling \(h>0\), and mass \(m\geq0\), define
\begin{equation}
\label{eq:deformed-charge}
\begin{aligned}
 D_v(h,m)&=\sum_i\theta(A_i(v))\cdot p_i
 +h\sum_{i<j}\theta(B_{ij}(v))\cdot[x_i,x_j]
 +m\sum_i\theta(L_i(v))\cdot x_i,\\
 \cH(h,m)&=\frac1{16}\sum_{\nu=1}^{16}D_{v_\nu}(h,m)^2.
\end{aligned}
\end{equation}
The second expression initially denotes a quadratic form on the
gauge-invariant smooth compactly supported core; \((v_\nu)\) is any
orthonormal real spin basis. Its closure will have the same notation.
The minus sign in \(B_{ij}\) is the transpose sign in
\((\gamma^{ij})^t=-\gamma^{ij}\), so \(D_v(1,0)\) is exactly the
linear combination of the charges in \eqref{eq:model-charge} with
coefficients \(v\).

For a real skew spatial matrix \(F\), use the rotation conventions
\begin{equation}
\label{alg:rotation-generator}
 S(F)=\frac14\sum_{i,j}F_{ij}\gamma^i\gamma^j,
 \qquad
 J(F)=p\cdot Fx+\frac{\ii}{2}\theta^t S(F)\theta.
\end{equation}
The color identity in the second term is implicit. In particular,
\(J\) includes both orbital and spin rotations. Gauge invariance means
equivariance of the Clifford-valued function, rather than invariance
of its scalar coordinate dependence alone. Infinitesimally this gives
\begin{equation}
\label{alg:gauge-generator}
 p\cdot\ad(w)x=-\frac{\ii}{2}\theta^t\ad(w)\theta
 \quad\text{on gauge-invariant functions},\qquad w\in\mathfrak k.
\end{equation}

\begin{proposition}[Square of a charge]
\label{prop:square}
On the gauge-invariant core,
\begin{equation}
\label{eq:square}
\begin{aligned}
 \cH(h,m)&=\frac12p^2+U_{h,m}
   +\frac{\ii}{2}\theta^t\left(\frac32mT-hK_j\gamma^j\right)\theta,
       \qquad K_j=\ad(x_j),\\
 2U_{h,m}&=\sum_{\mathrm{cyc}}
       \abs{h[x_a,x_b]-2mx_c}^2
       +h^2\sum_{\substack{i<j\\j>3}}\abs{[x_i,x_j]}^2
       +m^2\sum_{d=4}^9\abs{x_d}^2,\\
 D_v(h,m)^2&=\cH(h,m)+mJ(R(v)),
       \qquad R(v)_{ij}=s_jv^t\gamma^i\gamma^jTv.
\end{aligned}
\end{equation}
The cyclic sum uses \((a,b,c)=(1,2,3),(2,3,1),(3,1,2)\).
The matrix \(R(v)\) is skew, preserves \(\R^3\oplus\R^6\), and has
zero spin-basis average. The same formulas hold for an orthogonal
direct sum of compact color algebras and abelian color spaces.
\end{proposition}

\begin{proof}
\emph{The massless square.}
We first compute the kinetic and commutator contributions, including
the gauge cancellation. For a strictly increasing multi-index \(I\), let \(\gamma^I\) denote
the corresponding ordered gamma product. The products of ranks zero
through four form an orthogonal basis of the real \(16\)-by-\(16\)
matrices, for the trace scalar product. Indeed there are
\(1+9+36+84+126=256\) of them. Every nonconstant product of at most
eight distinct gammas is traceless: conjugation by a gamma outside the
product changes its sign for odd rank, and conjugation by a gamma
inside changes its sign for positive even rank. This also proves
orthogonality. The symmetric products have ranks zero, one, and four;
the skew products have ranks two and three.

Put \(\operatorname{sym}M=(M+M^t)/2\) and
\(\operatorname{skew}M=(M-M^t)/2\). The identity needed for the Yukawa
term is
\begin{equation}
\label{alg:spin-identity}
 \sum_i\operatorname{sym}\bigl((\gamma^iv)(\gamma^{ij}v)^t\bigr)
 =\frac12(v^t\gamma^jv)\Id-\frac12\gamma^j.
\end{equation}
To check it, test its left side against a symmetric monomial
\(M=\gamma^I\) of rank \(r\). Since
\(\sum_i\gamma^iM\gamma^i=(-1)^r(9-2r)M\), its trace pairing is
\[
 \bigl((-1)^r(9-2r)-1\bigr)v^tM\gamma^jv.
\]
For \(r=0,1,4\), this is respectively
\(8v^t\gamma^jv\), \(-8\delta_{Ij}\), and zero, which are precisely
the pairings with the right side of \eqref{alg:spin-identity}.

The CAR give \(\theta(w)^2=\abs{w}^2/2\). They therefore give the
kinetic term \(p^2/2\). The square of the bracket multiplier is
\(h^2\sum_{i<j}\abs{[x_i,x_j]}^2/2\): its possible remaining
four-form term is proportional to the alternating sum of
\(\ip{[x_i,x_j]}{[x_k,x_l]}\), and vanishes by invariance of the
color scalar product and the Jacobi identity.

The scalar, first-order part of the kinetic--bracket anticommutator is
\(h(v^t\gamma^jv)p\cdot K_jx\). In differentiating the bracket
coefficient with respect to \(x_i^A\), its component of color \(B\),
divided by \(h\), has spin coefficient
\(\sum_j(K_j)_{BA}\gamma^{ij}v\). Since \(K_j\) is skew in color,
its contraction with the kinetic Clifford multiplier selects the
symmetric spin matrix in \eqref{alg:spin-identity}. Consequently the
differentiated part is
\[
 -\frac{\ii h}{2}\theta^tK_j\gamma^j\theta
 +\frac{\ii h}{2}(v^t\gamma^jv)\theta^tK_j\theta.
\]
The second term cancels the first-order term by
\eqref{alg:gauge-generator}. This proves all the massless terms in
\eqref{eq:square}.

\smallskip\noindent\emph{The mass and rotation terms.}
The vectors \(L_i(v)\) have Gram matrix
\(\ip{L_i(v)}{L_j(v)}=s_i^2\delta_{ij}\). The bracket--mass
anticommutator is a scalar multiplier. Its color contractions are
alternating in three distinct spatial indices. Their spin coefficient
is proportional to
\[
 (s_i+s_j+s_k)v^t\gamma^i\gamma^j\gamma^kTv.
\]
If exactly one of the indices is among the first three, the sum of
the \(s\)'s is zero. If exactly two, or none, are among the first
three, the matrix in this expression is skew and its expectation is
zero. For \((i,j,k)=(1,2,3)\), \(T^2=-\Id\) gives
\(-6hm\ip{[x_1,x_2]}{x_3}\). Together with the bracket square and
the mass square, this is exactly \(U_{h,m}\) in
\eqref{eq:square}.

The scalar kinetic--mass cross term is \(mp\cdot R(v)x\). Diagonal
entries of \(R(v)\) vanish since \(T\) is skew. For mixed indices
in \(\R^3\) and \(\R^6\), \(\gamma^i\gamma^jT\) is skew; for
indices in the same block, \(s_i=s_j\), so antisymmetry in \(i,j\)
proves the stated form of \(R(v)\).
For the differentiated kinetic--mass term put
\[
 C(v)=\sum_i s_i\operatorname{skew}
               \bigl((\gamma^iv)(\gamma^iTv)^t\bigr).
\]
Its contribution is \(-\ii m\theta^tC(v)\theta\), and
\begin{equation}
\label{alg:mass-spin-identity}
 C(v)=-\frac34T-\frac12S(R(v)).
\end{equation}
Here is a trace verification that fixes both constants. For a skew
monomial \(M\) of rank \(r=2\) or \(3\), weighted conjugation,
using \(\sum_i s_i=0\), gives
\[
 \Tr(M^tC(v))
 =-2(-1)^r\left(\sum_{i\in I}s_i\right)v^tMTv.
\]
For rank three, only \(M=T\) can contribute: the case with one
first-three index has zero weight sum, and the other cases give a
skew product \(MT\). The pairing with \(T\) is \(-12\), giving
\(-3T/4\). For rank two, mixed blocks give zero, and in either pure
block the pairing is \(-4R(v)_{ij}\); since
\(\Tr((\gamma^{ij})^t\gamma^{ij})=16\) and
\(S(R)=\frac12\sum_{i<j}R_{ij}\gamma^{ij}\), this is the
\(-S(R)/2\) term. Thus the differentiated term is
\(3\ii m\theta^tT\theta/4+\ii m\theta^tS(R(v))\theta/2\), as
required.

Finally, \(\sum_\nu v_\nu v_\nu^t=\Id\). The trace of
\(\gamma^i\gamma^jT\) is zero for every \(i,j\), so the spin-basis
average of \(R(v_\nu)\) vanishes. Averaging the charge square gives
the asserted expression for \(\cH\). Only the invariant scalar
product and Jacobi identity were used in color space, proving the
last assertion as well.
\end{proof}

\subsection{The invariant charge and its coercive inequality}

The massless charges are covariant under \(\Spin(9)\); the massive
charges are covariant under \(\Spin(3)\times\Spin(6)\). We fix the
fermion grading as follows. Polarize the real spinor space by \(T\),
and realize the Clifford module as the exterior algebra on the
resulting eight-dimensional complex spinor space tensored with
\(\mathfrak k_\C\). Choose its empty vector to be the ground vector
of \(\ii\theta^tT\theta/2\). Gauge transformations and
\(\Spin(3)\times\Spin(6)\) preserve this polarization. Their actions
are the natural exterior actions: their connected semisimple groups
admit no character that could shift the infinitesimal spin action.
The gauge action also extends to \(\U(N)\) by conjugation; scalar
unitaries act trivially. Hence it factors through the adjoint gauge
group. The central minus one in \(\Spin(3)\) acts by minus the
identity on the one-particle space and therefore by exterior parity,
denoted \((-1)^F\).

To restrict a single massive charge to rotation invariants, we need a
subgroup fixing its spinor. This does not by itself remove the angular
term in its square. The choice below makes that term controllable by
averaging positive charge squares.

Write \(E_{12},E_{45},E_{67},E_{89}\) for the skew spatial matrices
whose indicated upper-triangular entry is one. The commuting complex
structures
\(\gamma^{12},\gamma^{45},\gamma^{67},\gamma^{89}\) have a common
real two-plane on which they agree. Indeed the three commuting
orthogonal projectors
\[
 \frac12(\Id-\gamma^{12}\gamma^{45}),\qquad
 \frac12(\Id-\gamma^{12}\gamma^{67}),\qquad
 \frac12(\Id-\gamma^{12}\gamma^{89})
\]
have product of trace two, by the gamma trace rules just proved.
Choose a real unit \(u\) in this plane. The nine-dimensional volume
element \(\gamma^1\cdots\gamma^9\) is a scalar \(\pm\Id\): it is
central, and the monomial basis shows that the commutant is scalar.
It follows that \(\gamma^3=\sigma\Id\) on the chosen plane, for
some \(\sigma\in\{1,-1\}\).

Let \(G\) be the connected subtorus in the spin group with Lie
algebra
\begin{equation}
\label{alg:fixing-torus}
 \left\{\sum_{q=1}^4t_qE_q:\ \sum_{q=1}^4t_q=0\right\},
 \qquad (E_1,E_2,E_3,E_4)=(E_{12},E_{45},E_{67},E_{89}).
\end{equation}
It fixes \(u\); covariance therefore implies that \(D_u(h,m)\)
commutes with \(G\). We always use this connected subgroup, with no
additional finite-component invariance imposed.

\begin{lemma}[Coercivity on rotation invariants]
\label{lem:coercive-charge}
On gauge- and \(G\)-invariant vectors, in the sense of quadratic
forms,
\begin{equation}
\label{eq:coercive-charge}
 \norm{D_u(h,m)\psi}^2\geq\frac23\cH(h,m)[\psi].
\end{equation}
\end{lemma}

\begin{proof}
The weight-plane relations give
\[
 R(u)=\sigma(-2E_{12}+E_{45}+E_{67}+E_{89}).
\]
Off-torus coefficients vanish by orthogonality of distinct weight
planes. On \(G\)-invariant vectors the generator of any coefficient
vector with zero sum is zero. Thus, putting
\(R_0=\sigma(E_{45}+E_{67}+E_{89})\),
\(J(R(u))=J(R_0)/3\). Average \(R(gu)\) over
\(g\in\Spin(3)\). The first block has average zero, and the
last-six block stays fixed, so the average is \(R_0\).
Proposition~\ref{prop:square} now gives on the invariant core
\[
 \int_{\Spin(3)}\norm{D_{gu}(h,m)\psi}^2\,\dd g
 =\cH(h,m)[\psi]+m\ip{\psi}{J(R_0)\psi}
 =3\norm{D_u(h,m)\psi}^2-2\cH(h,m)[\psi].
\]
Positivity proves the inequality. The domain argument below extends
it to the closed graph domain of the restricted charge.
\end{proof}

\subsection{Closed domains and the massless kernel}

For the undeformed full-space matrix model, see
Lundholm~\cite[Theorem~3.1]{Lundholm2010}. The constant-symbol cutoff
proof below also covers the massive charges and specifies the closed
realization used in the limiting arguments.

\begin{lemma}[Minimal and maximal domains]
\label{lem:domains}
For fixed \(h,m,v\), the charge \(D_v(h,m)\) is essentially
self-adjoint on the full smooth compactly supported core. Its closed
domain is its distributional maximal domain
\[
 \{\psi\in L^2:D_v(h,m)\psi\in L^2\text{ in distributions}\}.
\]
The same statement holds after restriction to any commuting compact
group of gauge or rotation symmetries. On gauge invariants at
\(m=0\), all unit-spinor charges have the same graph domain, and
\[
 \norm{D_v(h,0)\psi}^2=\cH(h,0)[\psi].
\]
Consequently the distributional \(L^2\) kernel of any one of these
charges is their common kernel. At \(h=1\) it is exactly \(\Ker H_N\)
for the form Hamiltonian in \eqref{eq:model-form}.
\end{lemma}

\begin{proof}
Write \(D=D_0+V(x)\), where \(V\) is a smooth self-adjoint matrix
and \(D_0\) has constant coefficients. Its symbol satisfies
\begin{equation}
\label{alg:symbol}
 \sigma_D(\xi)^2=\frac12\abs{\xi}^2\Id,
 \qquad
 [D,\chi]=-\ii\sum_{i,A}\theta^A(A_i(v))\partial_{i,A}\chi,
 \qquad
 \norm{[D,\chi](x)}_{\mathrm{op}}=\frac{\abs{\nabla\chi(x)}}{\sqrt2}
\end{equation}
for real scalar \(\chi\). If \(D\psi\in L^2\) distributionally,
localizing in a compact set makes \(V\psi\) square integrable. The
Fourier identity for \(D_0\), with a further compact cutoff, shows
that \(\psi\in H^1_{\mathrm{loc}}\).

Take radial cutoffs \(\chi_R\), equal to one on \(\abs{x}\leq R\)
and zero on \(\abs{x}\geq2R\), with
\(\abs{\nabla\chi_R}\leq C/R\). Then
\[
 D(\chi_R\psi)=\chi_RD\psi+[D,\chi_R]\psi
 \longrightarrow D\psi
\]
in \(L^2\), and \(\chi_R\psi\to\psi\) in \(L^2\). Each compactly
supported \(H^1\) vector can in turn be approximated in the graph
norm by convolution and a fixed outer cutoff; on that compact set
the potential is bounded and smooth. Thus the minimal domain equals
the maximal domain. The adjoint of the symmetric core operator has
precisely the distributional maximal domain, proving
self-adjointness. Averaging these approximants over a commuting
compact group preserves smooth compact support, and converges in
graph norm because its action commutes with \(D\).

At \(m=0\), Proposition~\ref{prop:square} makes the sixteen core
graph norms, and that of any unit \(v\), identical. Their common
completion is therefore both the closed single-charge graph domain
and the form domain of \(\cH(h,0)\). The associated nonnegative
operator has zero form precisely when every charge vanishes. This
proves the kernel statements. The same core approximation and
Lemma~\ref{lem:coercive-charge} imply that the restricted massive
graph domain is included in the averaged form domain and that
\eqref{eq:coercive-charge} holds there.
\end{proof}

The following is the explicit rotation primitive of
Hasler--Hoppe~\cite[Lemma~2(c)]{HaslerHoppe2002}, in the present
normalization. It supplies rotational invariance without any
assumption on \(\abs{x}\psi\). For a skew spatial matrix \(F\),
define the real spin matrices
\begin{equation}
\label{alg:explicit-P}
 P_i(F)=\frac1{16}\sum_jF_{ji}\gamma^j
        +\frac1{112}\sum_{j,k}F_{jk}\gamma^{ijk},
 \qquad
 M_v(F)=\sum_i\theta(P_i(F)v)\cdot x_i.
\end{equation}
Here \(\gamma^{ijk}\) is fully antisymmetrized, and is zero if two
indices coincide.

\begin{lemma}[An exact anticommutator for rotations]
\label{alg:rotation-exact}
For any orthonormal real spin basis,
\begin{equation}
\label{alg:rotation-exact-eq}
 \sum_{\nu=1}^{16}\{D_{v_\nu}(h,0),M_{v_\nu}(F)\}=J(F)
\end{equation}
on the smooth compactly supported core.
\end{lemma}

\begin{proof}
Orthogonality of gamma monomials gives
\[
 \Tr(\gamma^jP_i(F))=F_{ji},\qquad
 \Tr(\gamma^{jk}P_i(F))=0,\qquad
 \sum_i\gamma^iP_i(F)^t=-\frac12S(F).
\]
For the last identity, the rank-one term contributes
\(-S(F)/4\). For fixed distinct \(j,k\), exactly seven indices
\(i\) remain in the rank-three term, and
\(\gamma^i\gamma^{ijk}=\gamma^{jk}\); its transpose has a minus
sign. Hence this term contributes another \(-S(F)/4\).
The first trace identity yields the orbital term in the
anticommutator. The second kills its bracket multiplier. The
differentiated multiplier is
\(-\ii\theta^t(\sum_i\gamma^iP_i(F)^t)\theta\), which is exactly
the spin term of \(J(F)\).
\end{proof}

\begin{proposition}[Singlet property and parity]
\label{prop:singlet-parity}
Every gauge-invariant \(L^2\) vector in the undeformed kernel is
\(\Spin(9)\)-invariant and even. In particular it is \(G\)-invariant.
This assertion does not require finite-dimensionality of the kernel
or any moment of its vectors.
\end{proposition}

\begin{proof}
Let \(\mathcal Z\) be the common massless kernel furnished by
Lemma~\ref{lem:domains}. Covariance preserves \(\mathcal Z\).
Convolving its vectors with smooth approximate identities on the
compact group \(\Spin(9)\) gives a dense invariant subspace
\(\mathcal Z^\infty\) contained in the domains of all rotation
generators. For \(\phi,\psi\in\mathcal Z^\infty\), use the radial
cutoffs from Lemma~\ref{lem:domains} in
\eqref{alg:rotation-exact-eq}. Local regularity implies that
\(\chi_RM_v(F)\psi\) is in the charge domain. Pairing the
anticommutator with \(\phi\), and using that both vectors are killed
by every charge, gives
\[
 0=\ip{\phi}{\chi_RJ(F)\psi}
   +\sum_\nu\ip{\phi}{[D_{v_\nu}(h,0),\chi_R]M_{v_\nu}(F)\psi}.
\]
The multiplier in the last term has operator norm bounded by
\(C_F\abs{x}\abs{\nabla\chi_R}\), and is supported in
\(R\leq\abs{x}\leq2R\). Therefore that term is bounded in absolute
value by
\[
 C_F\norm{\mathbf1_{\{\abs{x}\geq R\}}\phi}
       \norm{\mathbf1_{\{\abs{x}\geq R\}}\psi}\longrightarrow0.
\]
Since \(J(F)\psi\in L^2\), the first term converges to
\(\ip{\phi}{J(F)\psi}\). The generator also preserves the closed
invariant subspace \(\mathcal Z\), so the vanishing of these matrix
elements for its dense smooth subspace implies \(J(F)\psi=0\).
All infinitesimal generators consequently vanish on the smooth
vectors of this unitary representation. Connectedness of
\(\Spin(9)\), followed by density, makes the representation trivial
on \(\mathcal Z\). The central minus one of \(\Spin(3)\) is exterior
parity, as established above; thus every vector of \(\mathcal Z\)
is even. The proof has used only escaping-annulus \(L^2\) tails of
the uncut states. This is the rotational-singlet argument of
Sethi--Stern~\cite{SethiStern2000}, in the explicit anticommutator form
of Hasler--Hoppe~\cite{HaslerHoppe2002}. The underlying cutoff principle
also appears in Hitchin's invariant $L^2$-cohomology argument
\cite[Theorem~3]{Hitchin2000}.
\end{proof}
